\scanpage{116}
% archive foliation: 115
Soit $X$ une ``courbe analytique'' i.e. espace $\mathbb C$-analytique \uline{régulier} partout de dimension 1. Soit
\[
t:X\longrightarrow\mathbb C
\]
un ``uniformisant'' i.e. \deleted{\unsure{une}} application analytique étale de $X$ dans $\mathbb C=\mathbb E^1(\mathbb C)$, i.e. une section $t$ de $\mathcal O_X$ telle que $dt\in\Gamma(X,\underbrace{\Omega^1_{X/\mathbb C}}_{\underline\omega})$ soit partout $\ne0$ i.e. soit une \uline{base} de $\Omega^1_{X/\mathbb C}$). Soit $s\in\Gamma(X,\mathcal O_X)$ -- \deleted{de préférence} un autre uniformisant, et posons
\[
G(t,s)=\left(-2\frac{s'''}{s'}+3\left(\frac{s''}{s'}\right)^2\right)dt^2\in\Gamma(X,\unsure{\underline\omega^{\otimes2}|\Omega^{\otimes2}})
\]
où $s'=\frac{ds}{dt}$, $s''=\frac{ds'}{dt}$, $s'''=\frac{ds''}{dt}$. On \deleted{\unsure{montre}} vérifie facilement, \deleted{que si $a,b$ sont des constantes, $a\ne0$} \inserted{si $a\in\mathbb C^*$, $b\in\mathbb C$}
\[
(1)\quad\left\{\begin{array}{ll}
G(t,\deleted{\unsure{as}}s+b)=G(t,s)&a)\\
G(t,as)=G(t,s)&b)\\
G(t,\frac1s)=G(t,s)&c)
\end{array}\right.
\]
où, dans la dernière équation, on suppose que $1/s$ est un uniformisant, i.e. $s$ n'a pas de zéros sur $X$. On vérifie de plus que, si $t,s,r$ sont trois uniformisants, on a
\[
(2)\qquad G(t,s)+G(s,r)=G(t,r)
\]
\unsure{Enfin} (3) la formation de $G(t,s)$ commute à la restriction à un ouvert. \deleted{\unsure{En}}

De (1) \inserted{c)} et (3) résulte qu'on peut prolonger l'expression $G(t,s)$ au cas où $s$, au lieu d'être un ``uniformisant'' i.e. un morphisme étale $X\longrightarrow\mathbb E^1(\mathbb C)$, est un morphisme étale
\[
X\longrightarrow\mathbb P^1(\mathbb C)
\]
i.e. une fonction méromorphe n'ayant que des pôles simples, et uniformisante en dehors de ceux-ci -- (dans \inserted{voisinage des} pôles de $s$, on pose $G(t,s)=G(t,1/s)$), et alors \inserted{\unsure{il ne forme plus de} \deleted{\unsure{fonction analytique}} \unsure{où $s$ est holomorphe})} (1) \deleted{devient} reste valable, et \uline{plus généralement}
\[
G\left(t,\frac{as+b}{cs+d}\right)=G(t,s)\qquad
\begin{gathered}a,b,c,d\in\mathbb C\\ad-bc\in\mathbb C^*.\end{gathered}
\]
